\documentclass[11pt,a4paper]{article}
\usepackage[left=24mm,right=24mm,top=22mm,bottom=20mm,includefoot]{geometry}
\usepackage{amsmath,amssymb,amsthm,booktabs}
\usepackage{iftex}
\ifPDFTeX
  \usepackage{CJKutf8}
  % CJKutf8 supplies its own Unicode map for each Chinese subfont.
  \makeatletter
  \AddToHook{cmd/CJK@addcmap/before}{\pdfnobuiltintounicode\font@name}
  \makeatother
\else
  \usepackage[no-math]{fontspec}
  \IfFontExistsTF{Noto Serif CJK SC}{
    \setmainfont{Noto Serif CJK SC}
    \setsansfont{Noto Sans CJK SC}
  }{
    \setmainfont{FandolSong-Regular.otf}[BoldFont=FandolSong-Bold.otf]
    \setsansfont{FandolHei-Regular.otf}[BoldFont=FandolHei-Bold.otf]
  }
  \ifXeTeX
    \XeTeXlinebreaklocale "zh"
    \XeTeXlinebreakskip=0pt plus 0.5pt
  \fi
\fi
\pagestyle{plain}
\raggedbottom
\linespread{1.13}
\setlength{\parindent}{2em}
\setlength{\parskip}{3pt plus 1pt}
\setlength{\jot}{3pt}
\DeclareMathSizes{10.95}{10.95}{9}{7}
% Use a single stroke for Chinese headings when pdfLaTeX has no bold CJK font.
\newcommand{\headingweight}{\ifPDFTeX\mdseries\else\bfseries\fi}
\makeatletter
\renewcommand{\section}{\@startsection{section}{1}{0pt}
  {2.0ex plus .3ex minus .2ex}{.9ex}{\normalfont\large\headingweight}}
\makeatother

\newcommand{\indep}{\mathrel{\perp\!\!\!\perp}}
\newtheoremstyle{report}{7pt}{7pt}{\normalfont}{}{\headingweight}{。}{0.4em}{}
\theoremstyle{report}
\newtheorem{assumption}{假设}
\newtheorem{theorem}{定理}
\newtheorem{lemma}{引理}
\newtheorem{proposition}{命题}
\newtheorem*{markova}{假设 3a}
\newtheorem*{markovb}{假设 3b}
\renewcommand{\proofname}{证明}
\makeatletter
\renewenvironment{proof}[1][证明]{\par\pushQED{\qed}\normalfont
  \topsep=6pt plus 2pt\relax\trivlist
  \item[\hskip\labelsep{\headingweight #1。}]\ignorespaces}
  {\popQED\endtrivlist\@endpefalse}
\makeatother
\newcommand{\pr}{\mathrm{pr}}
\newcommand{\E}{\mathrm E}
\newcommand{\avec}{\boldsymbol a}
\newcommand{\zero}{\boldsymbol 0}
\newcommand{\one}{\boldsymbol 1}
\begin{document}
\ifPDFTeX\begin{CJK*}{UTF8}{gbsn}\fi
\begin{center}
{\LARGE\headingweight 单路径 Markov 条件下的有序结局责任识别}\\[5pt]
{\large 从两个、三个暴露到一般 $p$ 个暴露}
\end{center}

结局等级按相邻一级计 $1$ 分。责任以该评分为单位；这一评分约定是归因定义的一部分。

\section{两个暴露：两种可选的单路径假设}
设 $X_1,X_2\in\{0,1\}$，$Y\in\{0,1,2\}$，数值越大表示结局越严重。
记 $T=(Y_{00},Y_{01},Y_{10},Y_{11})$，并定义
\[
\mathcal E_1=\{X_1=X_2=1,Y=1\},\qquad
\mathcal E_2=\{X_1=X_2=1,Y=2\}.
\]
对给定类型 $t$，将下式中的潜在结果取为该类型的相应分量：
\begin{equation}\label{eq:p2-score}
\begin{aligned}
R(1\mid t,\mathcal E_k)&=\tfrac12\{(Y_{10}-Y_{00})+(Y_{11}-Y_{01})\},\\
R(2\mid t,\mathcal E_k)&=\tfrac12\{(Y_{01}-Y_{00})+(Y_{11}-Y_{10})\},\\
R(0\mid t,\mathcal E_k)&=Y_{00},\qquad k=1,2,\quad t_{11}=k.
\end{aligned}
\end{equation}
其中 $0$ 表示背景。病例平均责任定义为
\[
\mathcal R_j(\mathcal E_k)
=\sum_{t:t_{11}=k}R(j\mid t,\mathcal E_k)\pr(T=t\mid\mathcal E_k),
\qquad j=0,1,2.
\]

\begin{assumption}[一致性、随机分配与正值性]\label{ass:p2-data}
观测结局为 $Y=Y_{X_1X_2}$，且
\[
(X_1,X_2)\indep T,\qquad
\pr(X_1=a,X_2=b)>0,\quad a,b=0,1.
\]
\end{assumption}

\begin{assumption}[单步单调增量]\label{ass:p2-unit}
下列不等式几乎处处成立：
\[
\begin{aligned}
0\le Y_{10}-Y_{00}\le1,\qquad&0\le Y_{11}-Y_{10}\le1,\\
0\le Y_{01}-Y_{00}\le1,\qquad&0\le Y_{11}-Y_{01}\le1.
\end{aligned}
\]
\end{assumption}

\begin{markova}[经 $Y_{10}$ 的单路径条件]
选定路径 $Y_{00}\longrightarrow Y_{10}\longrightarrow Y_{11}$，假设
\begin{equation}\label{eq:markov-a}
Y_{11}\indep Y_{00}\mid Y_{10}.
\end{equation}
\end{markova}

\begin{markovb}[经 $Y_{01}$ 的单路径条件]
选定路径 $Y_{00}\longrightarrow Y_{01}\longrightarrow Y_{11}$，假设
\begin{equation}\label{eq:markov-b}
Y_{11}\indep Y_{00}\mid Y_{01}.
\end{equation}
\end{markovb}
假设 3a 与假设 3b 是两种可替代的建模选择，识别时指定采用其中一种，
不要求同时成立。路径箭头表示暴露组合的加入顺序及相应的分布分解，
不表示一个潜在结果在因果上导致另一个潜在结果。
\setcounter{assumption}{3}

\clearpage
\section{两个暴露的识别定理}
记
\[
\xi_{ab}=\pr(Y_{ab}=2),\qquad
\theta=\pr(Y_{00}=0,Y_{11}=2).
\]
\begin{theorem}[选定单路径下的病例平均责任]\label{thm:p2}
在假设~\ref{ass:p2-data}、\ref{ass:p2-unit} 下，若采用假设 3a，则
\begin{equation}\label{eq:theta-a}
\theta=
\begin{cases}
\dfrac{\{\pr(Y_{00}=0)-\pr(Y_{10}=0)\}(\xi_{11}-\xi_{10})}
      {\pr(Y_{10}=1)},&\pr(Y_{10}=1)>0,\\[8pt]
0,&\pr(Y_{10}=1)=0.
\end{cases}
\end{equation}
若采用假设 3b，则
\begin{equation}\label{eq:theta-b}
\theta=
\begin{cases}
\dfrac{\{\pr(Y_{00}=0)-\pr(Y_{01}=0)\}(\xi_{11}-\xi_{01})}
      {\pr(Y_{01}=1)},&\pr(Y_{01}=1)>0,\\[8pt]
0,&\pr(Y_{01}=1)=0.
\end{cases}
\end{equation}
将所选模型给出的 $\theta$ 代入下式，即可识别全部平均责任。
当 $\xi_{11}>0$ 时，
\begin{equation}\label{eq:p2-grade2}
\begin{aligned}
\mathcal R_1(\mathcal E_2)
 &=\frac{\xi_{11}-\xi_{01}+\xi_{10}-\xi_{00}+\theta}{2\xi_{11}},\\
\mathcal R_2(\mathcal E_2)
 &=\frac{\xi_{11}-\xi_{10}+\xi_{01}-\xi_{00}+\theta}{2\xi_{11}},\\
\mathcal R_0(\mathcal E_2)
 &=\frac{\xi_{11}+\xi_{00}-\theta}{\xi_{11}}.
\end{aligned}
\end{equation}
当 $\pr(Y_{11}=1)>0$ 时，
\begin{equation}\label{eq:p2-grade1}
\begin{aligned}
\mathcal R_1(\mathcal E_1)
 &=\frac{\E(Y_{10})-\E(Y_{00})+\E(Y_{11})-\E(Y_{01})
          -\xi_{11}+\xi_{01}-\xi_{10}+\xi_{00}-\theta}
         {2\pr(Y_{11}=1)},\\[3pt]
\mathcal R_2(\mathcal E_1)
 &=\frac{\E(Y_{01})-\E(Y_{00})+\E(Y_{11})-\E(Y_{10})
          -\xi_{11}+\xi_{10}-\xi_{01}+\xi_{00}-\theta}
         {2\pr(Y_{11}=1)},\\[3pt]
\mathcal R_0(\mathcal E_1)
 &=\frac{\E(Y_{00})-\xi_{11}-\xi_{00}+\theta}{\pr(Y_{11}=1)}.
\end{aligned}
\end{equation}
每个病例等级下均有 $\sum_{j=0}^{2}\mathcal R_j(\mathcal E_k)=k$。
\end{theorem}

\begin{proof}
随机分配与一致性给出
$\pr(T=t\mid\mathcal E_k)=\pr(T=t\mid Y_{11}=k)$，
并识别四个潜在结果的全部边际概率和均值。
单步增量条件蕴含
\[
\{Y_{00}=0,Y_{11}=2\}
=\{Y_{00}=0,Y_{10}=1,Y_{11}=2\},
\]
且在该事件上同样有 $Y_{01}=1$。此外，
\[
\begin{aligned}
\pr(Y_{00}=0,Y_{10}=1)&=\pr(Y_{00}=0)-\pr(Y_{10}=0),\\
\pr(Y_{10}=1,Y_{11}=2)&=\xi_{11}-\xi_{10}.
\end{aligned}
\]
因此，在 $\pr(Y_{10}=1)>0$ 且采用假设 3a 时，
\[
\begin{aligned}
\theta
&=\pr(Y_{10}=1)\pr(Y_{00}=0,Y_{11}=2\mid Y_{10}=1)\\
&=\pr(Y_{10}=1)\pr(Y_{00}=0\mid Y_{10}=1)
                       \pr(Y_{11}=2\mid Y_{10}=1)\\
&=\frac{\pr(Y_{00}=0,Y_{10}=1)\pr(Y_{10}=1,Y_{11}=2)}
        {\pr(Y_{10}=1)}.
\end{aligned}
\]
这给出式~\eqref{eq:theta-a}。分母为零时，上述事件包含关系给出 $\theta=0$。
在假设 3b 下，将上述推导中的 $Y_{10}$ 换为 $Y_{01}$，得到式~\eqref{eq:theta-b}。

对 $ab=10,01$，单步增量条件给出
\[
\E\{Y_{ab}\mathbf1(Y_{11}=2)\}
=\pr(Y_{ab}=1,Y_{11}=2)+2\pr(Y_{ab}=2)
=\xi_{11}+\xi_{ab}.
\]
无暴露结局满足
\[
\begin{aligned}
\E\{Y_{00}\mathbf1(Y_{11}=2)\}
&=\pr(Y_{00}=1,Y_{11}=2)+2\xi_{00}\\
&=(\xi_{11}-\theta-\xi_{00})+2\xi_{00}
=\xi_{11}+\xi_{00}-\theta.
\end{aligned}
\]
将这些恒等式代入式~\eqref{eq:p2-score} 的带指标期望，再除以
$\pr(Y_{11}=2)$，得到式~\eqref{eq:p2-grade2}。

当 $Y_{11}=0$ 时，各责任均为零。
记式~\eqref{eq:p2-score} 中的任一责任为 $R_j(T)$，则
\[
\E\{R_j(T)\mathbf1(Y_{11}=1)\}
=\E\{R_j(T)\}-\E\{R_j(T)\mathbf1(Y_{11}=2)\}.
\]
将刚才计算的未标准化二级病例贡献相减，再除以 $\pr(Y_{11}=1)$，
得到式~\eqref{eq:p2-grade1}，该推导不要求 $\xi_{11}>0$。
逐类型有 $R_0(T)+R_1(T)+R_2(T)=Y_{11}$，取条件期望即得加总等式。
\end{proof}

两个 $\theta$ 表达式仅在同时施加假设 3a、3b 时才必须相等。
“选定一条路径并假设其成立”不同于“至少一条未知路径满足条件”；
后一模型需要对各候选路径及其相容联合分布分别分析，不能直接选取一个表达式作为点识别结果。

\section{三个暴露的病例平均责任}\label{sec:three}
设三个暴露依次为 $X_1,X_2,X_3$，结局取 $0,1,2$ 级。下标记录哪些暴露已加入。
例如，$Y_{101}$ 是仅有 $X_1$ 和 $X_3$ 时的潜在结果。固定 $k=1$ 或 $2$，记
\[
  \mathcal E_k=\{X_1=X_2=X_3=1,\ Y=k\},\qquad \pr(Y_{111}=k)>0.
\]
假设一致性成立，暴露分配独立于八个潜在结果，且每个组合的概率均为正。
因此，每个暴露组都给出相应潜在结果的边际分布。任一潜在结果在 $\mathcal E_k$ 下的均值，
等于它在 $Y_{111}=k$ 的个体中的均值。

\medskip
\noindent{\headingweight 三个暴露分别改变哪个潜在结果。}
下表列出加入各暴露时的全部增量。单步增量假设要求表中每个差都取 $0$ 或 $1$。
\begin{center}
\renewcommand{\arraystretch}{1.19}
\setlength{\tabcolsep}{5pt}
\begin{tabular}{@{}c c cc c@{}}
\toprule
暴露 & 最先加入 & \multicolumn{2}{c}{第二个加入} & 最后加入\\
\cmidrule(lr){3-4}
$X_1$ & $Y_{100}-Y_{000}$ & $Y_{110}-Y_{010}$ & $Y_{101}-Y_{001}$ & $Y_{111}-Y_{011}$\\
$X_2$ & $Y_{010}-Y_{000}$ & $Y_{110}-Y_{100}$ & $Y_{011}-Y_{001}$ & $Y_{111}-Y_{101}$\\
$X_3$ & $Y_{001}-Y_{000}$ & $Y_{101}-Y_{100}$ & $Y_{011}-Y_{010}$ & $Y_{111}-Y_{110}$\\
\bottomrule
\end{tabular}
\end{center}

要分配病例责任，需要知道表中各个潜在结果在 $Y_{111}=k$ 的个体中平均为多少。
终点 $Y_{111}$ 本身等于 $k$。$Y_{110},Y_{101},Y_{011}$ 距终点都只差一个暴露，
各自与 $Y_{111}$ 的联合分布由单步增量和两组边际分布确定。
例如，当 $Y_{111}=2$ 时，$Y_{110}$ 只能为 $1$ 或 $2$，而 $Y_{110}=2$ 必有 $Y_{111}=2$。
因此，对有二级病例的总体，
\[
  \pr(Y_{110}=1\mid Y_{111}=2)
  =\frac{\pr(Y_{111}=2)-\pr(Y_{110}=2)}{\pr(Y_{111}=2)}.
\]
还须识别与终点相隔两步的 $Y_{100},Y_{010},Y_{001}$，以及相隔三步的 $Y_{000}$，
在这组病例中的分布。这四个起点各对应下面一条条件。

\clearpage
\noindent{\large\headingweight 三个暴露的四条终点 Markov 条件}\par
固定优先顺序 $X_1,X_2,X_3$。起点尚无 $X_1$ 时就加入 $X_1$。
起点 $100$ 已经有 $X_1$，此时加入 $X_2$。例如，从无暴露到全暴露的选定路径为
\[
Y_{000}\xrightarrow{+X_1}Y_{100}\xrightarrow{+X_2}Y_{110}\xrightarrow{+X_3}Y_{111}.
\]
计算从终点往回进行。下表按计算次序列出对同一总体施加的四条条件。
\begin{center}
\renewcommand{\arraystretch}{1.20}
\setlength{\tabcolsep}{7pt}
\begin{tabular}{@{}cccl@{}}
\toprule
起点 & 加入 & 终点 Markov 条件 & 识别的二元分布\\
\midrule
$100$ & $X_2$ & $Y_{111}\indep Y_{100}\mid Y_{110}$ & $(Y_{100},Y_{111})$\\
$010$ & $X_1$ & $Y_{111}\indep Y_{010}\mid Y_{110}$ & $(Y_{010},Y_{111})$\\
$001$ & $X_1$ & $Y_{111}\indep Y_{001}\mid Y_{101}$ & $(Y_{001},Y_{111})$\\
\addlinespace[2pt]
$000$ & $X_1$ & $Y_{111}\indep Y_{000}\mid Y_{100}$ & $(Y_{000},Y_{111})$\\
\bottomrule
\end{tabular}
\end{center}

以上表第一行为例，在 $Y_{110}=s$ 的个体中，额外知道 $Y_{111}=k$ 后，
$Y_{100}$ 取各等级的概率仍为 $\pr(Y_{100}=r\mid Y_{110}=s)$。
在病例中按 $Y_{110}$ 的等级加权，即得第一行的识别式。四行分别给出
\[
\begin{aligned}
\pr(Y_{100}=r\mid Y_{111}=k)
 &=\sum_s \pr(Y_{100}=r\mid Y_{110}=s)\,\pr(Y_{110}=s\mid Y_{111}=k),\\
\pr(Y_{010}=r\mid Y_{111}=k)
 &=\sum_s \pr(Y_{010}=r\mid Y_{110}=s)\,\pr(Y_{110}=s\mid Y_{111}=k),\\
\pr(Y_{001}=r\mid Y_{111}=k)
 &=\sum_s \pr(Y_{001}=r\mid Y_{101}=s)\,\pr(Y_{101}=s\mid Y_{111}=k),\\
\pr(Y_{000}=r\mid Y_{111}=k)
 &=\sum_s \pr(Y_{000}=r\mid Y_{100}=s)\,\pr(Y_{100}=s\mid Y_{111}=k).
\end{aligned}
\]
这里 $r,s$ 取 $0,1,2$，求和时只计权重为正的等级。
前三式的右侧均由相邻暴露组合给出。最后一式用到的
$\pr(Y_{100}=s\mid Y_{111}=k)$ 已由第一式算出。

取第一式的 $r=0,\ k=2$。从 $Y_{100}=0$ 经 $Y_{110}$ 到 $Y_{111}=2$ 时，
中间等级只能是 $1$。在 $\pr(Y_{110}=1)>0$ 时，
\[
\begin{aligned}
\pr(Y_{100}=0\mid Y_{111}=2)
 &=\pr(Y_{100}=0\mid Y_{110}=1)\,\pr(Y_{110}=1\mid Y_{111}=2)\\
 &=\frac{\pr(Y_{100}=0)-\pr(Y_{110}=0)}{\pr(Y_{110}=1)}
   \times\frac{\pr(Y_{111}=2)-\pr(Y_{110}=2)}{\pr(Y_{111}=2)}.
\end{aligned}
\]
若 $\pr(Y_{110}=1)=0$，该病例概率为 $0$。
例如，若 $100,110,111$ 三组在 $0,1,2$ 级上的比例依次为
$(1/2,1/2,0)$、$(1/4,1/2,1/4)$ 和 $(1/8,3/8,1/2)$，
则上式的两个因子都等于 $1/2$。全暴露的二级病例中，仅保留 $X_1$ 时为 $0$ 级的比例为 $1/4$。

\medskip
\noindent{\headingweight 三个暴露各分多少分。}
固定上述病例等级，简记 $\mu_{abc}=\E(Y_{abc}\mid Y_{111}=k)$，其中 $\mu_{111}=k$。
六种加入顺序等权。一个暴露最先加入有两种顺序，最后加入也有两种。
在中间加入时，另外两个暴露各有一种顺序排在它前面。因此，增量表中四列的
权重依次为 $1/3,1/6,1/6,1/3$。逐行取病例中的平均值得
\[
\begin{aligned}
\mathcal R_1(\mathcal E_k)
 &=\tfrac13(\mu_{100}-\mu_{000})
   +\tfrac16\{(\mu_{110}-\mu_{010})+(\mu_{101}-\mu_{001})\}
   +\tfrac13(k-\mu_{011}),\\
\mathcal R_2(\mathcal E_k)
 &=\tfrac13(\mu_{010}-\mu_{000})
   +\tfrac16\{(\mu_{110}-\mu_{100})+(\mu_{011}-\mu_{001})\}
   +\tfrac13(k-\mu_{101}),\\
\mathcal R_3(\mathcal E_k)
 &=\tfrac13(\mu_{001}-\mu_{000})
   +\tfrac16\{(\mu_{101}-\mu_{100})+(\mu_{011}-\mu_{010})\}
   +\tfrac13(k-\mu_{110}),\\
\mathcal R_0(\mathcal E_k)&=\mu_{000}.
\end{aligned}
\]
每种顺序的三个增量都加总为 $Y_{111}-Y_{000}$。计入背景责任后，四项平均责任加总为 $k$。

\clearpage
\section{一般 $p$：暴露组合与责任}
设 $\boldsymbol X=(X_1,\ldots,X_p)\in\{0,1\}^p$，
潜在结果为 $Y_{\avec}\in\{0,1,\ldots,K\}$。
下标 $\avec=(a_1,\ldots,a_p)$ 直接表示暴露组合。
记 $|\avec|=\sum_j a_j$，$\zero=(0,\ldots,0)$，
$\one=(1,\ldots,1)$。当 $a_j=0$ 时，$\avec^{j+}$ 表示仅将第 $j$ 个坐标改为 $1$。
令 $T=(Y_{\avec}:\avec\in\{0,1\}^p)$，以及
\[
\mathcal E_k=\{\boldsymbol X=\one,Y=k\},\qquad k=1,\ldots,K.
\]
对 $j=1,\ldots,p$，定义
\begin{equation}\label{eq:general-score}
R_j(T)=
\sum_{\avec:a_j=0}
 \frac{|\avec|!\,(p-|\avec|-1)!}{p!}
 (Y_{\avec^{j+}}-Y_{\avec}),\qquad R_0(T)=Y_{\zero}.
\end{equation}
在 $Y_{\one}=k$ 的类型上，上式给出 $R(j\mid T,\mathcal E_k)$。
定义病例平均责任为
\[
\mathcal R_j(\mathcal E_k)=\E\{R_j(T)\mid\mathcal E_k\}.
\]

\begin{assumption}[一般 $p$ 的观测条件]\label{ass:p-data}
$Y=Y_{\boldsymbol X}$，$\boldsymbol X\indep T$，且
$\pr(\boldsymbol X=\avec)>0$ 对每个暴露组合成立。
\end{assumption}

\begin{assumption}[一般 $p$ 的单步增量]\label{ass:p-unit}
对每个 $a_j=0$ 的组合及暴露 $j$，几乎处处有
\[
0\le Y_{\avec^{j+}}-Y_{\avec}\le1.
\]
\end{assumption}

记
\[
f_{\avec}(r)=\pr(Y_{\avec}=r),\qquad
F_{\avec}(r)=\pr(Y_{\avec}\le r).
\]
在假设~\ref{ass:p-data} 下，
$f_{\avec}(r)=\pr(Y=r\mid\boldsymbol X=\avec)$，因而全部 $f_{\avec},F_{\avec}$ 可识别。

\begin{lemma}[相邻潜在结果的联合分布]\label{lem:edge}
若 $U,V\in\{0,\ldots,K\}$ 且 $0\le V-U\le1$，
则它们的边际分布识别联合分布。对 $0\le r<K$，
\begin{equation}\label{eq:edge}
\begin{aligned}
\pr(U=r,V=r+1)&=F_U(r)-F_V(r),\\
\pr(U=r,V=r)&=\pr(U=r)-\{F_U(r)-F_V(r)\}.
\end{aligned}
\end{equation}
另有 $\pr(U=K,V=K)=\pr(U=K)$，其余位置的概率为零。
\end{lemma}
\begin{proof}
单调性给出 $\{V\le r\}\subseteq\{U\le r\}$。
单步增量进一步给出
$\{U\le r,V>r\}=\{U=r,V=r+1\}$，所以第一式成立。
固定 $U=r$ 后只可能有 $V=r,r+1$，故第二式成立。
\end{proof}

\section{为每个组合指定一条通向全暴露的路径}
预先固定一个选择规则 $h$：对每个 $\avec\ne\one$，
$h(\avec)$ 只将 $\avec$ 的一个零坐标改为 $1$。
反复应用 $h$，每个组合均沿唯一的选定路径到达 $\one$。
这些选定边组成一棵指向全暴露组合的有向树。

\begin{assumption}[选定路径上的终点 Markov 条件]\label{ass:terminal}
对每个 $|\avec|\le p-2$ 的组合，假设
\begin{equation}\label{eq:terminal}
Y_{\one}\indep Y_{\avec}\mid Y_{h(\avec)}.
\end{equation}
全部条件均针对同一个潜在结果联合分布。
\end{assumption}
该条件要求：沿选定路径，下一节点的等级足以概括当前节点中与最终结局有关的信息。
第~\ref{sec:three}~节取 $h(100)=110$、$h(010)=110$、$h(001)=101$ 和 $h(000)=100$。
代入式~\eqref{eq:terminal}，恰好得到该节列出的四条条件。
每个起点只采用一个下一节点，不要求同一起点的所有加入方向均满足条件。
当 $|\avec|=p-1$ 时，$h(\avec)=\one$，上述条件独立关系自动成立，无需另加假设。

定义选定边的转移矩阵
\[
Q_{\avec}(r,s)=
\pr\{Y_{h(\avec)}=s\mid Y_{\avec}=r\},
\qquad \avec\ne\one.
\]
当 $f_{\avec}(r)>0$ 时，引理~\ref{lem:edge} 给出
\begin{equation}\label{eq:q}
\begin{aligned}
Q_{\avec}(r,r+1)
 &=\frac{F_{\avec}(r)-F_{h(\avec)}(r)}{f_{\avec}(r)}
       &&(r<K),\\
Q_{\avec}(r,r)
 &=1-Q_{\avec}(r,r+1)
       &&(r<K),\\
Q_{\avec}(K,K)&=1 &&(f_{\avec}(K)>0).
\end{aligned}
\end{equation}
其余位置为零。当 $f_{\avec}(r)=0$ 时，可将该行任取为一个概率向量，例如单位点质量。
此约定不影响下文出现的可识别联合概率。

\begin{theorem}[任意 $p$ 的病例 Shapley 责任]\label{thm:general}
在假设~\ref{ass:p-data}--\ref{ass:terminal} 下，定义
\begin{equation}\label{eq:h-recursion}
H_{\one}(r,k)=\mathbf1(r=k),\qquad
H_{\avec}(r,k)=
\sum_{s=0}^{K}Q_{\avec}(r,s)H_{h(\avec)}(s,k),
\quad \avec\ne\one.
\end{equation}
则对每个暴露组合及等级，
\begin{equation}\label{eq:pair-identification}
\pr(Y_{\avec}=r,Y_{\one}=k)=f_{\avec}(r)H_{\avec}(r,k).
\end{equation}
因此，当 $f_{\one}(k)>0$ 时，令
\begin{equation}\label{eq:mu}
\mu_{\avec}(k)=
\frac{\sum_{r=0}^{K}r f_{\avec}(r)H_{\avec}(r,k)}
     {f_{\one}(k)},
\end{equation}
即有
\begin{equation}\label{eq:general-r}
\begin{aligned}
\mathcal R_j(\mathcal E_k)
 &=\sum_{\avec:a_j=0}
   \frac{|\avec|!\,(p-|\avec|-1)!}{p!}
   \{\mu_{\avec^{j+}}(k)-\mu_{\avec}(k)\},
       &&j=1,\ldots,p,\\
\mathcal R_0(\mathcal E_k)&=\mu_{\zero}(k).
\end{aligned}
\end{equation}
所有右侧均由观测分布识别，且
$\sum_{j=0}^{p}\mathcal R_j(\mathcal E_k)=k$。
\end{theorem}

\begin{proof}
由引理~\ref{lem:edge}，每条选定边的联合分布可识别。
在 $f_{\avec}(r)>0$ 的行上，转移矩阵因而可识别。
对满足 $f_{h(\avec)}(s)=0$ 的等级 $s$，有
\[
\sum_{r=0}^{K}f_{\avec}(r)Q_{\avec}(r,s)
=f_{h(\avec)}(s)=0.
\]
非负性保证 $f_{\avec}(r)>0$ 时 $Q_{\avec}(r,s)=0$。
因此，正概率起点不会经一步转移到零概率等级，零概率行的约定不影响最终的联合概率。

对距全暴露组合仅一步的节点，
\[
H_{\avec}(r,k)=Q_{\avec}(r,k)
=\pr(Y_{\one}=k\mid Y_{\avec}=r).
\]
假设对节点 $h(\avec)$ 已证明
$H_{h(\avec)}(s,k)=\pr(Y_{\one}=k\mid Y_{h(\avec)}=s)$，
其中只需考虑正概率的 $s$。当 $|\avec|\le p-2$ 时，
全概率公式与假设~\ref{ass:terminal} 给出
\[
\begin{aligned}
&\pr(Y_{\one}=k\mid Y_{\avec}=r)\\
&\quad=\sum_{s=0}^{K}
 \pr(Y_{\one}=k\mid Y_{\avec}=r,Y_{h(\avec)}=s)
 \pr(Y_{h(\avec)}=s\mid Y_{\avec}=r)\\
&\quad=\sum_{s=0}^{K}
 \pr(Y_{\one}=k\mid Y_{h(\avec)}=s)Q_{\avec}(r,s)
 =H_{\avec}(r,k).
\end{aligned}
\]
各项仅在相应条件事件概率为正时使用条件概率；其余项的权重为零。
按节点距 $\one$ 的步数递推，得到全部节点的条件分布。
乘以 $f_{\avec}(r)$，即得式~\eqref{eq:pair-identification}，
并由求和得到 $\mu_{\avec}(k)=\E(Y_{\avec}\mid Y_{\one}=k)$。

随机分配与一致性保证
$\E\{R_j(T)\mid\mathcal E_k\}=\E\{R_j(T)\mid Y_{\one}=k\}$。
式~\eqref{eq:general-score} 对潜在结果是线性的，
故逐项取条件期望得到式~\eqref{eq:general-r}。
沿任一加入顺序，全部增量之和均为 $Y_{\one}-Y_{\zero}$；
对所有顺序平均后仍有该恒等式。
加入背景责任，再取条件期望即得责任加总为 $k$。
\end{proof}

定理~\ref{thm:general} 只需识别每个潜在结果与全暴露结局的二元联合分布。
它不要求识别完整类型分布，也未声称选定路径条件识别任意多个潜在结果的联合分布。

当 $p=2$ 时，唯一需要实质性条件独立假设的起点为 $00$。
选择 $h(00)=10$ 即得到假设 3a；选择 $h(00)=01$ 即得到假设 3b。
因此，定理~\ref{thm:p2} 是上述选定路径结构在三个结局等级下的显式表达。

\section{选定路径条件的相容性}
一般 $p$ 下，至少还有两个暴露尚未加入的起点共有 $2^p-p-1$ 个。
在每个这样的起点指定一个下一节点，便得到定理~\ref{thm:general} 所用的充分条件。

上述假设允许非退化的相容联合分布。例如，令 $B\in\{0,\ldots,K\}$
与 $D_1,\ldots,D_p\in\{0,1\}$ 相互独立，并定义
\[
Y_{\avec}=\min\left\{K,B+\sum_{j=1}^{p}a_jD_j\right\}.
\]
该构造满足所有单步增量限制。对任一选定边，尚未加入的 $D_j$
独立于 $(Y_{\avec},Y_{h(\avec)})$，且
\[
Y_{\one}=\min\left\{K,Y_{h(\avec)}+
               \sum_{j:h(\avec)_j=0}D_j\right\},
\]
故终点 Markov 条件成立。此构造用于说明假设的相容性，不是定理额外要求的模型。

\clearpage
\section{仅有一条完整路径的识别范围}
固定一种加入顺序，将该路径上的潜在结果依次记为
$Z_0,\ldots,Z_p$，其中 $Z_0=Y_{\zero}$、$Z_p=Y_{\one}$。
若该路径满足一阶 Markov 条件，则
\[
\pr(Z_0=z_0,\ldots,Z_p=z_p)
=\pr(Z_0=z_0)\prod_{m=1}^{p}
  \pr(Z_m=z_m\mid Z_{m-1}=z_{m-1}).
\]
单步增量与路径上各节点的边际分布识别每个相邻转移，
因而识别整条路径的联合分布，以及给定 $Z_p=k$ 时沿该固定顺序的每一步平均贡献。
一般 Shapley 责任还涉及该路径之外的暴露组合。

\begin{proposition}[一条完整 Markov 路径一般不足]\label{prop:one-path}
当 $p=3$、$Y\in\{0,1,2\}$ 时，即使所有单步增量均在 $[0,1]$ 内，
且指定路径
$Y_{000}\longrightarrow Y_{100}\longrightarrow Y_{110}\longrightarrow Y_{111}$
满足一阶 Markov 条件，四项病例平均 Shapley 责任也不一定全部可识别。
\end{proposition}
\begin{proof}
令
\[
Y_{000}=0,\qquad
Y_{100}=Y_{001}=Y_{110}=Y_{101}=Y_{011}=1.
\]
只保留 $Y_{010}=U\in\{0,1\}$ 与 $Y_{111}=V\in\{1,2\}$ 的随机性。
比较下列两个总体，各表格单元内的两种取值分别以概率 $1/2$ 出现：
\begin{center}
\renewcommand{\arraystretch}{1.25}
\begin{tabular}{@{}ccc@{}}
\toprule
总体 & $(U,V)$ 的取值 & 给定 $V=2$ 时的 $U$\\
\midrule
A & $(0,1),\ (1,2)$ & $1$\\
B & $(1,1),\ (0,2)$ & $0$\\
\bottomrule
\end{tabular}
\end{center}
两个总体均满足全部相邻单步增量限制，八个潜在结果的边际分布完全相同。
采用相同的独立随机暴露分配后，它们生成相同的观测分布。
指定路径在两个总体中均为 $(0,1,1,V)$，故均满足一阶 Markov 条件，
且整条路径的联合分布也相同。

在 $V=2$ 时，由 Shapley 权重 $1/3,1/6,1/3$ 直接计算，
\[
R_1=\frac{5-U}{6},\qquad
R_2=\frac{1+U}{3},\qquad
R_3=\frac{5-U}{6},\qquad R_0=0.
\]
例如，
\[
R_2=\tfrac13(U-0)+\tfrac16\{(1-1)+(1-1)\}+\tfrac13(2-1).
\]
因此，在 $\mathcal E_2=\{X_1=X_2=X_3=1,Y=2\}$ 下，
\[
\begin{aligned}
(\mathcal R_1,\mathcal R_2,\mathcal R_3)_{\mathrm A}
 &=\left(\tfrac23,\tfrac23,\tfrac23\right),\\
(\mathcal R_1,\mathcal R_2,\mathcal R_3)_{\mathrm B}
 &=\left(\tfrac56,\tfrac13,\tfrac56\right).
\end{aligned}
\]
观测分布及选定路径分布相同，而目标责任不同，故不能点识别。
\end{proof}

定理~\ref{thm:general} 中的“每个组合具有一条选定路径”覆盖全部需要归因的潜在结果。
命题~\ref{prop:one-path} 中的“一条完整路径”只覆盖 $p+1$ 个组合。
两者分别对应不同的信息要求，不能互相替代。
在 $p=2$、三级结局的情形中，一条路径能够补足唯一未识别的 $\theta$，
从而识别全部责任；一般 $p$ 不具备这一自动简化。

\ifPDFTeX\end{CJK*}\fi
\end{document}

